\chapter*{Innovations}
\addcontentsline{toc}{chapter}{Innovations}

The detector produces span-level, categorised hallucination flags with calibrated per-character probabilities for any VLM answer. It needs only the image, the question and the answer, not the model that wrote the answer. This makes it applicable in several settings:
\begin{itemize}
    \item \textbf{Quality assurance of VLM products.} Answers can be flagged for human review before they reach users, for example in accessibility applications that describe images to blind and low-vision users. There, a misread sign or a wrong count is a safety issue.
    \item \textbf{Data curation.} Hallucinated spans can be filtered or corrected in VLM-generated captions and instruction data before they are used to train other models. Category labels show which ability (OCR, counting, object grounding) needs more data.
    \item \textbf{Model evaluation.} The rate and type of hallucinations can be measured per language and per category for a candidate model, without manual annotation of every answer.
\end{itemize}
The components run on a single 24\,GB GPU, and the vote and the probe add no training beyond the judge and a small probe, which keeps deployment cost moderate.
